\chapter[Exact Agent-Network Renormalization]{Exact Renormalization of\\Agent Networks}
\label{ch:agent-network-rg}

This chapter constructs an effective variational theory for the interaction
model of \Cref{ch:local-collective-elbo}.  Its primary objects are normalized
measures and Markov kernels.  Actions, couplings, and differential beta
functions appear only after a reference and a scale identification have been
declared.  Exact closure is obtained by retaining induced hyperedges, marked
attention events, complete holonomy data, and path-space memory.  A pairwise,
memoryless model with one scalar attention row and one averaged group link is
therefore a truncation, not the general effective theory.

\section{Law-level coarse-graining and the VFE identity}
\label{sec:rg-law-level-vfe}

Let $P(do,dy)$ be the fixed normalized joint law on standard-Borel spaces, let
$o$ be regular, let $\Pi_o(dy)$ be its posterior, and let
$Q_o\ll\Pi_o$ be a recognition law.  A coarse channel between standard-Borel
latent spaces is a fixed Markov kernel
\begin{equation}
C:\mathsf Y\rightsquigarrow\mathsf Z
\label{eq:rg-coarse-channel}
\end{equation}
that does not read $Q_o$ and does not alter the observation coordinate.  Define
\begin{equation}
P^c(do,dz)=\int_{\mathsf Y}C(y,dz)P(do,dy),
\qquad
\Pi_o^c=\Pi_oC,
\qquad
Q_o^c=Q_oC.
\label{eq:rg-common-pushforward}
\end{equation}

\theoremheading{Exact coarse VFE and discarded conditional information}{thm:rg-exact-coarse-vfe}
Let the fine and coarse VFE expressions be defined as extended-real sums, with
the common finite evidence term separated from their KL gaps, and let
\begin{equation}
\widehat Q_o(dy,dz)=Q_o(dy)C(y,dz),
\qquad
\widehat\Pi_o(dy,dz)=\Pi_o(dy)C(y,dz).
\label{eq:rg-bridge-lifts}
\end{equation}
Then $P^c$ has the same observation evidence as $P$, and the extended-real
additive identity is
\begin{equation}
\Fenergy_P(Q_o)
=\Fenergy_{P^c}(Q_o^c)+\int_{\mathsf Z}
\KL\bigl(\widehat Q_o(dy\given z)
\Vert\widehat\Pi_o(dy\given z)\bigr)Q_o^c(dz)
.
\label{eq:rg-vfe-chain-rule}
\end{equation}
In particular, the fine VFE is at least the coarse VFE.  If the fine KL is
finite, their ordinary real-valued difference equals the displayed
conditional KL.
The corresponding coarse ELBO rises only because a conditional inference gap
has been discarded; the evidence has not increased.  \status{ESTABLISHED}

\paragraph{Proof.}
The observation marginal is unchanged because $C(y,\mathsf Z)=1$.  Attaching
the same channel to both measures preserves the fine relative entropy,
$\KL(Q_o\Vert\Pi_o)=\KL(\widehat Q_o\Vert\widehat\Pi_o)$.  Disintegration on
$z$ and the relative-entropy chain rule split the latter into
$\KL(Q_o^c\Vert\Pi_o^c)$ plus the integral in
\eqref{eq:rg-vfe-chain-rule}.  Substitute the exact evidence identities at the
two scales.  $\square$

The theorem does not cover different generative and recognition channels, a
fitted rather than pushed-forward $P^c$, or a simultaneous coarse-graining of
the observation.  Such operations generally change the comparison target or
the evidence event.  An invertible or evidence-sufficient observation channel
is an exception, but its preserved event and posterior must be proved
explicitly.  \status{ESTABLISHED}

\section{The exact effective likelihood and action}
\label{sec:rg-effective-action}

Let $\rho=P_0$ be a probability reference and let
\begin{equation}
m_o(dy)=L_o(y)\rho(dy),
\qquad
L_o(y)=\exp[-H_o(y)]
\label{eq:rg-success-submeasure}
\end{equation}
be the unnormalized evidence submeasure.  Its mass is $Z(o)$.  Push both
measures through $C$:
\begin{equation}
\rho^c=\rho C,
\qquad
m_o^c=m_oC.
\label{eq:rg-reference-evidence-pair}
\end{equation}
Absolute continuity $m_o^c\ll\rho^c$ follows from
$m_o\ll\rho$.  Define
\begin{equation}
L_o^c(z)=\frac{dm_o^c}{d\rho^c}(z),
\qquad
H_o^c(z)=-\log L_o^c(z).
\label{eq:rg-effective-likelihood-action}
\end{equation}

\theoremheading{Conditional-partition formula and composition}{thm:rg-effective-action}
Let $R_\rho(z,dy)$ be a reverse conditional of $Y$ given $Z=z$ under
$\rho(dy)C(y,dz)$.  Then
\begin{equation}
\exp[-H_o^c(z)]
=\int\exp[-H_o(y)]R_\rho(z,dy)
\label{eq:rg-conditional-partition}
\end{equation}
$\rho^c$-almost everywhere, and
\begin{equation}
Z(o)=\int\exp[-H_o^c(z)]\rho^c(dz).
\label{eq:rg-effective-evidence}
\end{equation}
For a second coarse channel $D$, the direct and staged effective likelihoods
agree almost everywhere when their conditional versions are chosen from the
same joint tower.  \status{ESTABLISHED}

\paragraph{Proof.}
For bounded measurable $f$,
\[
\int f(z)m_o^c(dz)
=\iint f(z)L_o(y)C(y,dz)\rho(dy)
=\int f(z)\E_\rho[L_o(Y)\given Z=z]\rho^c(dz).
\]
This identifies the Radon--Nikodym derivative and proves
\eqref{eq:rg-conditional-partition}.  Setting $f=1$ proves
\eqref{eq:rg-effective-evidence}.  The tower property proves composition.
$\square$

For binary interaction successes, $0\leq L_o\leq1$ implies
$0\leq L_o^c\leq1$, so one effective binary interaction record realizes the
entire coarse action.  For general record densities, $L_o^c$ remains a valid
density relative to the pushed probability reference, although it need not be
bounded by one.  An infinite Lebesgue reference cannot be pushed through an
arbitrary projection without checking sigma-finiteness; this is why the
probability pair $(\rho,m_o)$ is primary.  \status{ESTABLISHED}

\section{Exact closure generates hyperedges}
\label{sec:rg-hypergraph-closure}

Suppose the complete nonnegative density of the law, including the baseline,
is factorized against a product or block-product reference.  Marginalizing an
internal variable multiplies every incident factor, including baseline
factors, and integrates their product.  The result is one factor on the union
of the remaining incident scopes.  Iterating this operation represents the
normalized law by generally unnormalized factors on the elimination cliques;
the inherited total normalizer, not each factor, normalizes their product.
Tonelli's theorem makes the resulting joint law independent of elimination
order even though the displayed factorization is scheme dependent.  For an
arbitrary correlated baseline
without such a factorization, retain the baseline as one global factor and use
the law-level pushforward instead.  \status{ESTABLISHED}

For a finite meta-agent set $\mathcal P$ and a finite-valued effective action,
choose an anchor $z^\circ$ and define
\begin{equation}
\Phi_A^c(z_A)
=\sum_{B\subseteq A}(-1)^{|A|-|B|}
H_o^c(z_B,z_{B^c}^\circ),
\qquad A\subseteq\mathcal P.
\label{eq:rg-mobius-potentials}
\end{equation}
Mobius inversion gives
\begin{equation}
H_o^c(z)=\sum_{A\subseteq\mathcal P}\Phi_A^c(z_A).
\label{eq:rg-full-hypergraph-action}
\end{equation}
Thus the full hypergraph family is exactly closed.  If $H_o^c$ is extended
valued, retaining it as one top-order factor gives exact closure without
forming undefined alternating sums.  \status{ESTABLISHED}

The anchor is a covariant background, not a coordinatewise constant under
frame changes.  Consequently each displayed factor obeys
$\Phi_A^c(gz_A;gz^\circ)=\Phi_A^c(z_A;z^\circ)$ when the action is
gauge invariant.  If an anchor is held fixed only in one coordinate chart,
gauge covariance is asserted only for the reconstructed sum
\eqref{eq:rg-full-hypergraph-action}.  \status{ESTABLISHED}

Pairwise closure is false in general.  Eliminating the center $s_0$ of an
Ising star produces
\begin{equation}
\sum_{s_0=\pm1}\exp\left[s_0\left(h_0+\sum_{r=1}^3J_rs_r\right)\right]
=2\cosh\left(h_0+\sum_{r=1}^3J_rs_r\right),
\label{eq:rg-hidden-star-factor}
\end{equation}
whose negative log has cubic coefficient
$2\operatorname{sech}^2(h_0)\tanh(h_0)J_1J_2J_3+O(J^5)$ when
$h_0J_1J_2J_3\ne0$; it is therefore nonzero for all sufficiently small
nonzero $J_r$ with $h_0\ne0$.  \status{ESTABLISHED}

\section{Fine--meta bridge kernels}
\label{sec:rg-cross-scale-kernels}

The posterior bridge
\begin{equation}
B_o(dy,dz)=\Pi_o(dy)C(y,dz)
\label{eq:rg-posterior-bridge}
\end{equation}
defines the full reverse kernel $R_o(z,dy)=B_o(dy\given z)$.  For fine agent
$j$ and meta-agent $I$, take the $(Y_j,Z_I)$ marginal and disintegrate it in
both directions:
\begin{equation}
\Pi_{o,j}(dy_j)K_{I\leftarrow j}(y_j,dz_I)
=\Pi_{o,I}^c(dz_I)K_{j\leftarrow I}(z_I,dy_j).
\label{eq:rg-agent-meta-bayes-adjoint}
\end{equation}
These normalized kernels are the exact agent-to-meta and meta-to-agent
pair-marginal association, or bridge, kernels.  They are not thereby causal
generative interactions.  For bounded observables $f$ and $g$,
they obey
\begin{equation}
\int f(y_j)(K_{I\leftarrow j}g)(y_j)\Pi_{o,j}(dy_j)
=\int (K_{j\leftarrow I}f)(z_I)g(z_I)\Pi_{o,I}^c(dz_I).
\label{eq:rg-agent-meta-adjointness}
\end{equation}
\status{ESTABLISHED}

Equation~\eqref{eq:rg-agent-meta-bayes-adjoint} does not say that $Z_I$ alone
reconstructs the full fine posterior.  The exact full reverse generally reads
all meta variables.  A product baseline and blockwise coarse channel factorize
the baseline reverse, but a likelihood coupling blocks can destroy that
factorization in the posterior.  Posterior-local refinement additionally
requires the conditional-independence property
$Y_I\perp Z_{-I}\mid(Z_I,o)$, or another separately proved sufficient
condition.  \status{ESTABLISHED}

For a deterministic statistic $z=c(y)$, the observable pullback and posterior
conditional expectation are
\begin{equation}
(\mathsf P f)(y)=f(c(y)),
\qquad
(\mathsf Cg)(z)=\E_{\Pi_o}[g(Y)\given c(Y)=z].
\label{eq:rg-conditional-operators}
\end{equation}
They satisfy $\mathsf C\mathsf P=I$ in $L^p(\Pi_o^c)$, equivalently
$\Pi_o^c$-almost everywhere, and compose by the tower property along a fixed
nested filtration.  Values of a conditional kernel at unused coarse points
are immaterial.  For a randomized channel, the corresponding
operators remain adjoint but $\mathsf C\mathsf P=I$ is not automatic.
\status{ESTABLISHED}

\section{Gauge-covariant cross-scale operators}
\label{sec:rg-gauge-cross-scale}

For an arbitrary block partition $\mathcal P$, define the component-indexed
meta-agent set
\begin{equation}
\widehat{\mathcal P}
=\{\widehat I=(I,c):I\in\mathcal P,
c\in\pi_0(\Gamma_I)\}.
\label{eq:rg-component-meta-index}
\end{equation}
Choose a rooted spanning tree inside each connected component $\widehat I$ and
require successive trees to be nested.  From this point through the rooted
operator formulas, $I,J$ denote elements of $\widehat{\mathcal P}$ and
$r_I$ denotes their single component root.  The state of an original
disconnected block is the product over its component meta-agents; collapsing
that product to one root, row, or feature is an additional declared coarse
channel.  For channel $x\in\{b,m\}$, let
$\tau^x_{I\leftarrow i}$ be the ordered product of the channel-specific graph
links from $i$ to the root $r_I$.  Under the passive section rechoices
$a_i^b=a_i$ and $a_i^m=b_i$ of \Cref{eq:geo-local-reframing}, it transforms as
\begin{equation}
\tau^{x\prime}_{I\leftarrow i}
=(a^x_{r_I})^{-1}\tau^x_{I\leftarrow i}a_i^x.
\label{eq:rg-tree-transport-law}
\end{equation}
These are two passive coordinate rechoices in one principal $G$-bundle, not
an active $G\times G$ gauge symmetry.  For one active automorphism, the two
local functions instead obey $k_i^m=h_i^{-1}k_i^bh_i$.

Every non-tree edge is retained through the channel-typed based holonomy map
\begin{equation}
H_I^x:\pi_1(\Gamma_I,r_I)\longrightarrow G,
\qquad
H_I^{x\prime}(\gamma)
=(a^x_{r_I})^{-1}H_I^x(\gamma)a^x_{r_I}.
\label{eq:rg-full-holonomy-representation}
\end{equation}
For each marked boundary edge $e:j\to i$ from block $J$ to block $I$, retain
\begin{equation}
V_e^x=\tau^x_{I\leftarrow i}\Theta_e^x
(\tau^x_{J\leftarrow j})^{-1},
\qquad
V_e^{x\prime}=(a^x_{r_I})^{-1}V_e^xa^x_{r_J}.
\label{eq:rg-dressed-boundary-edge}
\end{equation}
The exact datum is the simultaneous root-framed equivariant state
$(\bar z_I^x,H_I^x,\{V_e^x\})$, or its one simultaneous root-gauge orbit;
separately quotienting each holonomy loses its orientation relative to the root
features and boundary legs.  For noncompact $G$, a naive conjugacy quotient
need not be standard Borel, so the measure tier retains the raw root-framed
$G$-space unless a proper-action or standard-Borel quotient theorem is
supplied.  Different trees change the presentation by Nielsen transformations;
nested trees with the compatibility below give strict composition, while
arbitrary choices give a root-gauge-equivariant isomorphism of presentations.
\status{ESTABLISHED}

On a separately declared linear feature fiber $V_x$, let
$R_x:G\to\operatorname{GL}(V_x)$ be a linear representation.  Choose
$w_{Ii}\geq0$ with $\sum_{i\in I}w_{Ii}=1$ and set
\begin{equation}
(\mathsf C_xz)_I
=\sum_{i\in I}w_{Ii}R_x(\tau^x_{I\leftarrow i})z_i,
\qquad
(\mathsf P_x\bar z)_i
=R_x(\tau^x_{I\leftarrow i})^{-1}\bar z_I.
\label{eq:rg-linear-cross-scale-maps}
\end{equation}
Then $\mathsf C_x\mathsf P_x=I$.  With
$R_{x,f}=\bigoplus_iR_x(a_i^x)$ and
$R_{x,c}=\bigoplus_IR_x(a_{r_I}^x)$, passive covariance is
\begin{equation}
\mathsf C_x'=R_{x,c}^{-1}\mathsf C_xR_{x,f},
\qquad
\mathsf P_x'=R_{x,f}^{-1}\mathsf P_xR_{x,c}.
\label{eq:rg-linear-cross-scale-covariance}
\end{equation}
At nested levels $0\to1\to2$, let
$\sigma^{x,12}_{A\leftarrow I}$ be the ordered interroot link from the root
of $I$ to the root of $A$.  Strict composition additionally requires
\begin{equation}
\tau^{x,02}_{A\leftarrow i}
=\sigma^{x,12}_{A\leftarrow I}\tau^{x,01}_{I\leftarrow i},
\qquad
w^{02}_{Ai}=w^{12}_{AI}w^{01}_{Ii}.
\label{eq:rg-linear-nested-compatibility}
\end{equation}
It then follows that
$\mathsf C_x^{02}=\mathsf C_x^{12}\mathsf C_x^{01}$ and
$\mathsf P_x^{02}=\mathsf P_x^{01}\mathsf P_x^{12}$; nested forests alone
would not imply either equality.  If $A_x$ is a fine message operator with
$A_x'=R_{x,f}^{-1}A_xR_{x,f}$, the
fine-to-meta, meta-to-fine, and meta-to-meta message operators are
\begin{equation}
A^x_{I\leftarrow j}=(\mathsf C_xA_x)_{Ij},
\qquad
A^x_{j\leftarrow I}=(A_x\mathsf P_x)_{jI},
\qquad
A^x_{I\leftarrow J}=(\mathsf C_xA_x\mathsf P_x)_{IJ}.
\label{eq:rg-cross-scale-message-operators}
\end{equation}
They transform between the corresponding endpoint frames; in particular,
\[
(\mathsf C_xA_x\mathsf P_x)'
=R_{x,c}^{-1}(\mathsf C_xA_x\mathsf P_x)R_{x,c}.
\]
\status{ESTABLISHED}

These message formulas are not an energy marginalization theorem.  A hard-tied
quadratic energy uses $\mathsf P_x^*H_x\mathsf P_x$; an exact Gaussian marginal
uses a Schur complement; a general law uses
\Cref{thm:rg-effective-action}.  The identity
$\mathsf C_x\mathsf P_x=I$ alone does not make
$\mathsf C_xH_x\mathsf P_x$ positive.  \status{ESTABLISHED}

If $\Phi_f:V_{b,f}\to V_{m,f}$ is a separately declared linear
cross-channel morphism on this tier, usually $\bigoplus_i\Phi_i$, with
$\Phi_f'=R_{m,f}^{-1}\Phi_fR_{b,f}$, its compressed map is
\begin{equation}
\Phi^c=\mathsf C_m\Phi_f\mathsf P_b,
\qquad
\Phi^{c\prime}=R_{m,c}^{-1}\Phi^cR_{b,c}.
\label{eq:rg-coarse-cross-channel-map}
\end{equation}
The common principal bundle does not create $\Phi$, and the two representations
must never be identified merely because both use $G$.  \status{ESTABLISHED}

\section{Exact attention between meta-agents}
\label{sec:rg-meta-attention}

A row-stochastic matrix $\beta_{ij}$ is a conditional law of a source given a
receiver.  It cannot be coarse-grained without the receiver law.  Let
$\alpha_i(y)$ be a normalized receiver occupancy and define the joint marked
edge-event law
\begin{equation}
\eta_{ij}(y)=\alpha_i(y)\beta_{ij}(y),
\qquad
\sum_{i,j}\eta_{ij}(y)=1.
\label{eq:rg-attention-event-law}
\end{equation}
The event probabilities are gauge-invariant scalars:
$\eta_{ij}'=\eta_{ij}$.
For a posterior bridge and node partition, set
\begin{equation}
\eta^c_{IJ}(z)
=\E_{B_o}\left[
\sum_{i\in I}\sum_{j\in J}\eta_{ij}(Y)
\middle|Z=z\right],
\quad
\alpha_I^c(z)=\sum_J\eta^c_{IJ}(z),
\quad
\beta^c_{IJ}(z)=\frac{\eta^c_{IJ}(z)}{\alpha_I^c(z)}.
\label{eq:rg-meta-attention}
\end{equation}
The last ratio is used on $\{\alpha_I^c>0\}$; a row on a zero-occupancy
receiver is immaterial.  Pushing $\eta$ and then disintegrating is normalized
and associative under nested partitions by the tower property.  The next
scale must push the joint law $\eta^c$, not $\beta^c$ alone.  Consequently
$\eta^c$, $\alpha^c$, and $\beta^c$ are gauge invariant.  \status{ESTABLISHED}

In an augmented receiver--source-label version of
\Cref{sec:local-attention-elbo}, the fixed recognition-independent label
baseline is
\begin{equation}
P_0^{\rm att}(dy,di,dj)=P_0^Y(dy)a_i\pi_{ij},
\qquad
\sum_i a_i=1,
\quad
\sum_j\pi_{ij}=1.
\label{eq:rg-attention-fixed-label-joint}
\end{equation}
Fix the latent state $y$ and define the row partition function
$Z_i(y)=\sum_j\pi_{ij}\exp[-D_{ij}(y)/\tau]$.  At the selected record, the
conditional prior inside a block with $a_I=\sum_{k\in I}a_k>0$ is
$a_{i\mid I}=a_i/a_I$.  Define the conditional block evidence weight
$W_I(y)=\sum_{k\in I}a_{k\mid I}c_k(o_k,y)Z_k(y)$.  On
$\{W_I(y)>0\}$, its posterior is evidence weighted:
\begin{equation}
\alpha_i'(y)=\frac{a_{i\mid I} c_i(o_i,y)Z_i(y)}
{W_I(y)},
\qquad
\beta^c_{IJ}(y)=\sum_{i\in I}\alpha_i'(y)
\sum_{j\in J}\beta_{ij}.
\label{eq:rg-attention-evidence-weights}
\end{equation}
The formula holds posterior-almost everywhere.  On a zero-evidence block the
receiver row is an arbitrary conditional version and is removed from the
effective support.
The equality $\alpha_i'=a_{i\mid I}$ holds exactly when
$c_i(o_i,y)Z_i(y)$ is constant on the conditional support in $I$.  This
guarantees equality with the pre-observation coarse
mixture, although equality can also occur accidentally after source
coarsening---for example, when all receivers induce the same block-source row.
If all receiver labels are retained, the exact meta-label is
their vector or multiset; replacing it by one categorical row is a further
coarse channel whose conditional KL appears in
\eqref{eq:rg-vfe-chain-rule}.  \status{ESTABLISHED}

For one categorical row and a source partition $j\mapsto J$, the exact coarse
prior and energy are
\begin{equation}
\pi^c_J=\sum_{j\in J}\pi_j,
\qquad
E^c_J=-\tau\log\left[
\frac{1}{\pi^c_J}\sum_{j\in J}\pi_j\exp[-E_j/\tau]
\right].
\label{eq:rg-attention-log-sum-exp}
\end{equation}
for every block with $\pi_J^c>0$; a zero-prior block has zero posterior mass and
is removed from the effective support.  This log-sum-exp rule associates for nested partitions of the same row.  It
does not merge independently normalized receiver rows without
\eqref{eq:rg-attention-evidence-weights}.  \status{ESTABLISHED}

The transported meta-message in channel $x$ is represented exactly by marked
parallel edges.  Let the fine mark $U^x_{i\leftarrow j}$ obey the passive
endpoint law and define its dressed version
$V_{ij}^x=\tau^x_{I\leftarrow i}U^x_{i\leftarrow j}
(\tau^x_{J\leftarrow j})^{-1}$.  On the linear feature tier, define the
following moment only under the conditional Bochner-integrability condition
\begin{equation}
\E_{B_o}\left[
\sum_{i\in I,j\in J}\eta_{ij}(Y)\|R_x(V_{ij}^x)\|
\middle|Z\right]<\infty
\quad\text{almost everywhere}.
\label{eq:rg-attention-bochner-domain}
\end{equation}
Then
\begin{equation}
\begin{split}
\mathsf T^x_{IJ}(z)
&=\E_{B_o}\left[
\sum_{i\in I,j\in J}\eta_{ij}(Y)
R_x(V_{ij}^x)\middle|Z=z\right],\\
\overline{\mathsf U}^x_{IJ}(z)
&=\frac{\mathsf T^x_{IJ}(z)}{\eta^c_{IJ}(z)}
\quad\text{on }\{\eta^c_{IJ}>0\},\\
\mathsf A^{x,\mathrm{att}}_{IJ}(z)
&=\frac{\mathsf T^x_{IJ}(z)}{\alpha_I^c(z)}
\quad\text{on }\{\alpha_I^c>0\}.
\end{split}
\label{eq:rg-attention-hom-operator}
\end{equation}
On $\{\eta^c_{IJ}>0\}$,
$\mathsf A^{x,\mathrm{att}}_{IJ}
=\beta^c_{IJ}\overline{\mathsf U}^x_{IJ}$.  On
$\{\alpha_I^c>0,\eta^c_{IJ}=0\}$, the attention-weighted block is zero and
$\overline{\mathsf U}_{IJ}^x$ is left undefined as a null-event conditional.
For $z'=R_{x,c}^{-1}z$, covariance is pointwise:
$\mathsf T_{IJ}^{x\prime}(z')=R_x(a_{r_I}^x)^{-1}\mathsf T_{IJ}^x(z)
R_x(a_{r_J}^x)$, so these operators transform from the $J$ root frame to
the $I$ root frame.  The conditional mean need not belong to $R_x(G)$: the
average of rotations by $\pi/2$ and $-\pi/2$ is the zero matrix.  Moreover,
this Hom operator is exact only for its stated first operator moment on a
common source input.  If marks and source features are correlated, their means
do not determine $\E[Uz]$; for example, equal-probability marks/features
$(I,v)$ and $(-I,-v)$ have zero separate means but mean product $v$.  Full
closure therefore retains the marked operator--feature kernel or all marked
edges, with \eqref{eq:rg-attention-hom-operator} as a derived moment, rather
than a fictitious averaged group element.  \status{ESTABLISHED}

At a fixed outside context, define
$\beta_I^Q(J\given z_{-I})=Q(J\given z_{-I})$ and
$K_{I\leftarrow J}^Q(dz_I\given z_{-I})
=Q(dz_I\given J,z_{-I})$, and define the corresponding generative
conditionals under $P(\cdot\given z_{-I},o)$.  When the recognition
categorical and conditional kernels are absolutely continuous with respect to
their generative counterparts, with the extended-value convention, the
conditional KL chain rule is
\begin{equation}
\begin{aligned}
&\KL\bigl(Q(J,Z_I\given z_{-I})\Vert
P(J,Z_I\given z_{-I},o)\bigr)\\
&\quad=\KL\bigl(\beta_I^Q(\cdot\given z_{-I})\Vert
\beta_I^P(\cdot\given z_{-I},o)\bigr)\\
&\qquad+\sum_J\beta_I^Q(J\given z_{-I})
\KL\bigl(K_{I\leftarrow J}^Q(\cdot\given z_{-I})\Vert
K_{I\leftarrow J}^P(\cdot\given z_{-I},o)\bigr).
\end{aligned}
\label{eq:rg-meta-attention-kl-split}
\end{equation}
Thus meta-attention and meta-agent interaction kernels are two typed parts of
one conditional ELBO, not interchangeable names for the same object.
\status{ESTABLISHED}

\begin{figure}[htbp]
\centering
\begin{tikzpicture}[>=stealth]
  \node[draw,rounded corners,fill=lightblue,minimum width=2.7cm,minimum height=0.8cm]
    (fine) at (0,2.1) {fine law $(\rho,m_o,\eta)$};
  \node[draw,rounded corners,fill=lightgreen,minimum width=2.7cm,minimum height=0.8cm]
    (coarse) at (5.2,2.1) {coarse law $(\rho^c,m_o^c,\eta^c)$};
  \node[draw,rounded corners,fill=lightorange,minimum width=2.7cm,minimum height=0.8cm]
    (action) at (5.2,0) {$H_o^c=-\log(dm_o^c/d\rho^c)$};
  \node[draw,rounded corners,fill=lightblue,minimum width=2.7cm,minimum height=0.8cm]
    (fixed) at (0,0) {rescaled common space};
  \draw[->,very thick,primaryblue] (fine)--node[above]{$C_\#$}(coarse);
  \draw[->,very thick,sciencegreen] (coarse)--node[right]{Radon--Nikodym}(action);
  \draw[->,very thick,cautionorange] (action)--node[midway,yshift=-0.32cm,fill=white,inner sep=1pt]{identification $I_b$}(fixed);
  \draw[->,very thick,primaryblue] (fixed)--node[left]{next exact step}(fine);
  \node[align=center,text width=5.2cm] at (2.6,-1.65)
    {Hyperedges, holonomy, joint attention events, and path memory remain in the state.};
\end{tikzpicture}
\caption{The exact RG acts first on normalized measures.  An effective action
and a beta function are introduced only after a pushed reference and a
declared identification return the changing scale spaces to one comparison
space.}
\label{fig:rg-exact-measure-flow}
\end{figure}

\section{Ordered-path closure and memory}
\label{sec:rg-path-space}

For a finite ordered path interval, apply the pointwise coarse channel to the
full joint path law.  Its order parameter is auxiliary: it is neither a
coordinate on the fixed base $\mathcal C$ nor the intrinsic information clock
of \Cref{ch:relational-inference}.  \Cref{thm:rg-exact-coarse-vfe} then applies
on path space without change.  The resulting coarse path law is normalized
and gauge covariant when the required spatial transports and, for a separately
declared dynamical model, parameter-direction transports are included.  It
need not be first-order Markov.  \status{ESTABLISHED}

For a deterministic state partition $c$ and a fine transition kernel $T$, a
coarse first-order kernel $T^c$ exists for every initial law exactly under
strong lumpability:
\begin{equation}
T(y,c^{-1}B)=T(y',c^{-1}B)
\quad\text{whenever }c(y)=c(y')
\label{eq:rg-strong-lumpability}
\end{equation}
for every coarse measurable set $B$.  Otherwise the exact effective action is
not first-order Markov for at least one initial law.  Particular weakly
lumpable initial laws can still yield a first-order coarse kernel even when
strong lumpability fails.  \status{ESTABLISHED}

On a linear observable tier with resolved projection $\Pi=\mathsf P\mathsf C$
and $\mathsf Q=I-\Pi$, elimination of the unresolved recurrence produces the
memory operators
\begin{equation}
\mathsf CT\mathsf Q(\mathsf QT\mathsf Q)^n
\mathsf QT\mathsf P,
\qquad n\geq0,
\label{eq:rg-memory-operators}
\end{equation}
together with terms depending on the unresolved initial state.  Autonomous
resolved dynamics on an admitted initial class holds exactly when all the
relevant memory and noise operators vanish there.  On the resolved initial
class $\mathsf Qx_0=0$, the stronger condition
$\mathsf QT\mathsf P=0$ is sufficient but not necessary.  For a general admitted
initial class one must additionally require
\begin{equation}
\mathsf CT\mathsf Q(\mathsf QT\mathsf Q)^n\mathsf Qx_0=0,
\qquad n\geq0,
\label{eq:rg-initial-noise-vanishing}
\end{equation}
throughout that class.  \status{ESTABLISHED}

\section{The RG transformation and beta functions}
\label{sec:rg-beta-function}

Let $b>1$ label one blocking ratio.  A genuine RG step consists of a coarse
channel $C_b$ and a declared rescaling/identification kernel $I_b$ that returns
the target to a common measurable state space.  Define the typed rescaled
kernel $K_b=C_bI_b$.  Compatibility means, after the declared canonical
identifications,
\begin{equation}
K_{b_1b_2}=K_{b_1}K_{b_2}.
\label{eq:rg-kernel-semigroup}
\end{equation}
On the primary measure pair,
define
\begin{equation}
\mathcal R_b(\rho,m_o)
=\bigl((\rho C_b)I_b,(m_oC_b)I_b\bigr)
=(\rho K_b,m_oK_b).
\label{eq:rg-measure-pair-map}
\end{equation}
Write $\mathcal R_b^\rho\rho=(\rho C_b)I_b=\rho K_b$ for its reference
component.
Equation~\eqref{eq:rg-kernel-semigroup}, together with compatible nested
forests, makes
$\mathcal R_{b_1b_2}=\mathcal R_{b_2}\mathcal R_{b_1}$; otherwise the sequence
is a typed cocycle rather than an autonomous semigroup.  \status{ESTABLISHED}

On a dominated tier, the action map retains its reference argument:
\begin{equation}
\mathcal R_b^H[H;\rho]
:=-\log\frac{d\bigl((e^{-H}\rho)K_b\bigr)}
{d(\rho K_b)}.
\label{eq:rg-reference-dependent-action-map}
\end{equation}
On a declared vector space of finite-valued actions that is closed under this
map for the declared reference trajectory and on which subtraction exists, the
exact discrete step beta functional is
\begin{equation}
\mathfrak B_b^H[H;\rho]
=\frac{\mathcal R_b^H[H;\rho]-H}{\log b}.
\label{eq:rg-discrete-beta-functional}
\end{equation}
If the evidence mass is deliberately forgotten and finite unnormalized actions
are quotiented by additive constants, the same formula is interpreted in that
quotient.  For extended-valued actions, including points at which the
numerator would be $+\infty-(+\infty)$, only the measure-pair RG is primary.  A beta
function on a truncated coupling vector exists only after invariance of that
finite family is proved; otherwise omitted hyperedges contribute an explicit
residual.  \status{ESTABLISHED}

Suppose instead that scale is a genuine continuous semigroup and that, relative
to one common dominating measure, the positive densities
$r_t,p_t\in\operatorname{Dom}(\mathcal A^*)$ are differentiable in a topology
that permits pointwise ratios and the following dual pairings, with
$p_t=r_t\exp[-H_t]$, $\dot r_t=\mathcal A^*r_t$, and
$\dot p_t=\mathcal A^*p_t$.  Then
\begin{equation}
\boxed{
\dot H_t
=-\frac{\mathcal A^*(r_t\exp[-H_t])}
{r_t\exp[-H_t]}
+\frac{\mathcal A^*r_t}{r_t}}
\label{eq:rg-continuous-beta-functional}
\end{equation}
almost everywhere wherever the displayed ratios and their dual pairings are
defined.  If the reference is stationary, the second term vanishes.
\status{ESTABLISHED}

\paragraph{Proof.}
Differentiate $H_t=-\log(p_t/r_t)$ and substitute the two forward equations.
$\square$

On a declared Banach action space with a Schauder basis
$\{\psi_A\}$ and continuous coordinate functionals
$\{\psi_A^*\}$, or a Hilbert space with a Riesz basis, suppose the expansion
$H_t=\sum_Ag_A(t)\psi_A$ converges and differentiation and coefficient
extraction can be interchanged.  The component beta functions are then the
coefficients of the right side of \eqref{eq:rg-continuous-beta-functional}:
\begin{equation}
\dot g_A=\left\langle\psi_A^*,
-\frac{\mathcal A^*(r_te^{-H_t})}{r_te^{-H_t}}
+\frac{\mathcal A^*r_t}{r_t}\right\rangle.
\label{eq:rg-component-beta}
\end{equation}
This is exact for the stated complete convergent basis.  A truncated basis
defines a projection and leaves an explicit residual.  \status{ESTABLISHED}

Attention flows primarily through the joint event law.  Along a differentiable
path on $\{\alpha_I>0\}$, since
$\beta_{IJ}=\eta_{IJ}/\alpha_I$,
\begin{equation}
\dot\beta_{IJ}
=\frac{\dot\eta_{IJ}-\beta_{IJ}\dot\alpha_I}{\alpha_I}.
\label{eq:rg-attention-beta-event}
\end{equation}
For a differentiable single softmax row on a fixed strictly positive prior
support, with $\tau>0$ and
$\beta_J=\pi_J\exp[-E_J/\tau]/Z$, this becomes the replicator beta function
\begin{equation}
\dot\beta_J=\beta_J\left(u_J-\sum_K\beta_Ku_K\right),
\qquad
u_J=\partial_t\log\pi_J-\frac{\dot E_J}{\tau}
+\frac{E_J\dot\tau}{\tau^2}.
\label{eq:rg-attention-beta-softmax}
\end{equation}
At a support change or on $\{\alpha_I=0\}$, the joint law $\eta$ is the
well-defined flow variable; neither the quotient row nor its logarithmic
parametrization is used.
\status{ESTABLISHED}

\section{Fixed points, flows, and scaling operators}
\label{sec:rg-fixed-points}

\theoremheading{Exhaustive fixed-point equations}{thm:rg-fixed-point-equations}
A normalized fixed theory is exactly a rescaled invariant pair
\begin{equation}
\mathcal R_b(\rho_*,m_*)=(\rho_*,m_*)
\quad\text{for every declared }b.
\label{eq:rg-fixed-measure-pair}
\end{equation}
On the dominated normalized tier this is equivalent to the conjunction
\begin{equation}
\mathcal R_b^\rho\rho_*=\rho_*,
\qquad
\mathcal R_b^H[H_*;\rho_*]=H_*,
\label{eq:rg-fixed-action}
\end{equation}
almost everywhere with respect to the fixed reference.  If the evidence mass
is discarded and actions are considered only projectively, the weaker
fixed-ray equation is
\begin{equation}
\mathcal R_b^H[H_*;\rho_*]=H_*+c_b.
\label{eq:rg-fixed-action-ray}
\end{equation}
It represents the normalized pair in \eqref{eq:rg-fixed-measure-pair} only when
the reference is also invariant and normalization fixes $c_b=0$.  A fixed
attention law is a fixed rescaled event measure $\eta_*$.  A fixed conditional
row alone does not make the attention law fixed: $\eta$ is fixed exactly when
its receiver marginal $\alpha$ is fixed and its disintegrated row $\beta$ is
fixed $\alpha$-almost everywhere.  \status{ESTABLISHED}

\paragraph{Proof.}
The first equation is the definition of invariance on the common rescaled
measure space.  Invariance of its first component and taking the
Radon--Nikodym derivative of its second component give the two conditions in
\eqref{eq:rg-fixed-action}.  Quotienting by the action's additive gauge gives
\eqref{eq:rg-fixed-action-ray}, but restoring the tracked evidence mass removes
that freedom.  Disintegration of a fixed $\eta$ gives the final statement.
$\square$

At a fixed pair, additionally require $H_*$ to lie in the declared
finite-valued action Banach space and
\begin{equation}
0<L_*^c(z)=\E_{\rho_*}[e^{-H_*}\given Z=z]<\infty
\quad \rho_*^c\text{-almost everywhere}.
\label{eq:rg-linearization-positive-likelihood}
\end{equation}
For bounded perturbations $\varphi\in L^\infty$ (or, more
generally, perturbations with conditional exponential integrability in a
neighborhood of zero and the first two required moments), the first two
variations of one exact conditional-partition step are
\begin{equation}
D_H\mathcal R_b^H\big|_{(H_*,\rho_*)}[\varphi](z)
=\E_{\Pi_*(dy\given z)}\varphi(Y),
\qquad
D_H^2\mathcal R_b^H\big|_{(H_*,\rho_*)}[\varphi,\varphi](z)
=-\operatorname{Var}_{\Pi_*(dy\given z)}\varphi(Y),
\label{eq:rg-linearized-action}
\end{equation}
followed by the declared rescaling.  On an infinite-dimensional Banach or
Hilbert action space, growth is classified by the full spectrum and spectral
radius of this derivative, including continuous and residual spectrum.  For an
isolated eigenoperator satisfying $D_H\mathcal R_b^H[\psi_a]
=\lambda_a\psi_a$, the exponent $y_a=\log|\lambda_a|/\log b$ is relevant,
marginal, or irrelevant according as it is positive, zero, or negative; the
sign or complex phase is retained.  Generalized eigenspaces are used only in
finite dimensions or for isolated point spectrum with the requisite spectral
projections.
\status{ESTABLISHED}

The equations characterize all fixed points without claiming that every model
class admits a closed-form enumeration.  The following exact sectors illustrate
them.  The identity channel fixes every law.  A terminal channel followed by
its one-point identification has the constant infrared fixed theory.  For
independent additive vector states and the block statistic
$b^{-1/\alpha}\sum_{i=1}^bY_i$, every strictly $\alpha$-stable baseline
with constant likelihood $m_o=Z(o)\rho$ gives a fixed measure pair; the
centered Gaussian is the $\alpha=2$ case.  At one fixed integer $b$ this
list is not exhaustive because $b$-semistable laws can also be fixed.
Exhaustiveness by strict stability requires invariance over all declared block
sizes (or a continuum of scales) with the usual centering and regularity
hypotheses.  \status{ESTABLISHED}

Flat graph connection data form an invariant sector under rooted contraction,
but become a fixed connection only after a self-similar graph identification
has been declared.  Uniform attention is fixed-form only when receiver
occupancies, block sizes, source labels, and self-edge conventions are also
self-similar.  One-hot rows are invariant boundary sectors of
\eqref{eq:rg-attention-beta-softmax}.  These qualifications prevent an
invariant form from being mislabeled as a complete fixed theory.
\status{ESTABLISHED}

Along every exact coarse path, \eqref{eq:rg-vfe-chain-rule} makes the recognition
gap nonincreasing and the ELBO nondecreasing at fixed evidence.  This monotone
flow is information loss under resolution, not a proof of approach to a
nontrivial critical fixed point.  Attraction requires spectral control of
\eqref{eq:rg-linearized-action} on the declared common space.
\status{ESTABLISHED}

\section[Closure theorem]{Closure theorem and infinite extensions}
\label{sec:rg-closure-theorem}

\theoremheading{Complete finite law-level gauge-VFE theory}{thm:rg-complete-effective-theory}
For every finite standard-Borel agent hypergraph, suppose the model supplies
the common principal bundle and its two associated statistical bundles;
normalized gauge-covariant baseline and interaction-record kernels with
positive finite evidence; channel-specific incidence and boundary links
$U^x,\Theta^x$ obeying
$U_{a\leftarrow i}^{x\prime}=(a_a^x)^{-1}U_{a\leftarrow i}^xa_i^x$
and the separate belief and model representations; normalized recognition laws
satisfying the stated absolute-continuity and log-integrability conditions;
structural coarse channels that are recognition independent and gauge
equivariant; globally gauge-equivariant jointly measurable versions of every
disintegration used for the displayed bridge kernels; a nested
component-rooted forest filtration; pushed probability references; the
simultaneous raw root-framed holonomy maps, root features,
and dressed boundary generators; joint marked attention-event laws; every
induced effective hyperedge; full path laws or their exact memory kernels; and
rescaling kernels satisfying
\eqref{eq:rg-kernel-semigroup}.  Then the global and local ELBOs,
agent--meta bridge kernels, meta-agent interaction kernels, meta-attention,
effective likelihoods, measure-pair scale composition, and measure-pair fixed
equations defined in this and the preceding chapter are simultaneously
normalized, gauge covariant, and exact.  \status{ESTABLISHED}

\paragraph{Proof.}
Normalization and the collective/local identities are
\Cref{prop:obs-interaction-normalization,thm:obs-collective-vfe,thm:obs-local-global-potential}.
The common-pushforward chain rule is \Cref{thm:rg-exact-coarse-vfe}; the
conditional-partition and tower identities are \Cref{thm:rg-effective-action}.
Mobius inversion or the retained top-order factor closes the hypergraph.
Regular disintegration gives an almost-everywhere version family of bridge
kernels, and the supplied equivariant-version hypothesis selects the displayed
globally covariant versions.  Rooted dressing gives the
separate gauge-covariant channel maps while the full holonomy representation
preserves the information lost by a tree presentation.  Pushforward and
disintegration of $\eta$ give exact attention.  Path pushforward supplies
dynamic closure, and the explicit common-space kernel identity types and
composes the measure-pair RG.  Each operation is closed in the stated
collection, so their finite composition is closed.  $\square$

\corollaryheading{Exact analytic action, beta, and scaling tier}{cor:rg-complete-analytic-tier}
In addition to the hypotheses of \Cref{thm:rg-complete-effective-theory},
suppose the pushed measures admit the declared common dominators.  On the
linear cross-scale tier, additionally supply the vector features and linear
representations $R_x$, multiplicative weights and inter-root transports in
\eqref{eq:rg-linear-nested-compatibility}, the extra fine cross morphism
$\Phi_f$, and the full marked operator--feature kernels.  These hypotheses
make the fine--meta, meta--fine, meta--meta, cross-channel, and attention
operators exact and compositional on their stated domains; the Hom moments
also require \eqref{eq:rg-attention-bochner-domain}.  For the discrete beta
require a declared reference trajectory and a finite-valued action vector
space closed under the reference-dependent action map.  For the continuous beta require the generator-domain,
differentiability, positivity, and dual-pairing hypotheses stated before
\eqref{eq:rg-continuous-beta-functional}.  For component beta functions
require the stated Schauder or Riesz basis and interchange conditions.  For
action fixed points retain invariance of the reference component.  For
linearization require \eqref{eq:rg-linearization-positive-likelihood} and use
the bounded or conditionally exponentially integrable perturbations stated
before \eqref{eq:rg-linearized-action}.  Then the
effective actions, discrete and continuous beta functionals, component flows,
action fixed-point equations, and linearized scaling operators are exact on
their declared analytic domains.  \status{ESTABLISHED}

\paragraph{Proof.}
The Radon--Nikodym formula gives the action.  Ordinary subtraction on the
finite action space gives the discrete beta; differentiation in the declared
generator domain gives the continuous beta.  Valid coefficient extraction
gives the component flow.  Invariance of both the reference and likelihood
ratio is equivalent to pair invariance, and dominated differentiation under
the conditional partition gives the first two variations.  $\square$

The theorem is uniform in the number and arrangement of agents; it is not a
verification on one fixed graph size.  For a countable agent or time index,
suppose instead that the normalized finite-cylinder joints form a projectively
consistent family on standard Borel coordinate spaces.  Kolmogorov extension
then supplies the full normalized joint \citep{Kallenberg2021}, and the stated
identities hold for every regular finite-cylinder observation and finite block.
Projective consistency alone does not make one prescribed infinite record
positive probability: an infinite sequence of independent positive-cost
binary successes can still be null.  Conditioning on an infinite record
therefore requires a chosen regular conditional at an almost-everywhere record
or a proved DLR, Doob-transform, or finite-volume posterior limit.  An
infinite-law KL pushforward additionally requires compatible
$Q\ll\Pi$ and an infinite coarse channel.  Existence and uniqueness of a
thermodynamic DLR state, convergence of free-energy densities, and interchange
of volume and RG limits remain the separate open obligations recorded in
\Cref{app:claim-ledger}.  \status{ESTABLISHED}
